\documentclass[11pt]{article}
\usepackage{amsmath,amssymb,hyperref}
\usepackage{geometry}
\usepackage{bbm}
\geometry{margin=1in}

% Generated by wolfbook-docs-agent from the live notebook on 2026-07-15.
% Reflects notebook state at generation time -- regenerate after material
% changes rather than hand-editing this file. Content corresponds to the
% companion XXX_CG_L2_Clean.md, but is restructured into continuous,
% topic-organized prose rather than a per-symbol transliteration, and every
% Greek-letter variable name from the notebook (lambda, theta, kappa, chi,
% tau) is rendered as a proper LaTeX math macro rather than a literal
% Unicode glyph. This file compiles cleanly with plain pdflatex.

\title{XXX\_CG\_L2\_Clean --- Reference}
\author{}
\date{Generated by wolfbook-docs-agent from the live notebook on 2026-07-15}

\newcommand{\qth}{Q_\theta}

\begin{document}
\maketitle

\section*{Overview}

This is the first Clean notebook in the XXX subproject: it sets up the two-site ($L=2$) rational
$\mathfrak{gl}(2)$ (XXX) spin chain --- independent representations $[\lambda_1,0]$, $[\lambda_2,0]$
at each site, independent inhomogeneities $\theta_1,\theta_2$, and a shared companion twist $G$ ---
and then solves the associated Baxter TQ equation for that chain. It was assembled on 2026-07-15 by
the \texttt{wolfbook-builder} agent from two \texttt{Experiments/} notebooks, kept unchanged as the
historical derivation record:

\begin{itemize}
\item \textbf{Section A of the notebook} (cells 2--26) --- ported verbatim from
  \texttt{Experiments/su2\_XXX\_L2.wb}: GT-basis/generator machinery (reused from
  \texttt{Experiments/XXX\_su2.wb}), Lax/monodromy/transfer matrices, the fusion (Hirota) hierarchy,
  the SoV $B$-operator, and the SoV pseudovacuum/covector basis, each followed by its own inline
  validation cell.
\item \textbf{Section B of the notebook} (cells 27--34) --- the seven routine-defining cells (no
  validation/regression cells) from \texttt{Experiments/Baxter\_Solver\_L2.wb}: the Baxter-equation
  ingredients, the genuine transfer-matrix-eigenstate solve (\texttt{TauEigensystem}), and the
  linear TQ solve for the Q-function (\texttt{QSolve}/\texttt{Q1}). Section B calls Section A's
  transfer matrix directly rather than re-deriving it (the duplicate transfer-matrix machinery that
  \texttt{Baxter\_Solver\_L2.wb} had ported verbatim was dropped during promotion, per
  \texttt{HANDOFF.md}).
\end{itemize}

Every construction in Section A that has a direct $L=1$ counterpart is cross-referenced against
\texttt{paper/xxxCG.tex} \S2 (``The XXX spin chain and separation of variables'') below; the paper
currently only states the $L=1$ case, so these cross-references are generalizations by tensor
product, not literal transcriptions. Section B (the Baxter/TQ solver) has no paper counterpart at
all yet --- \texttt{paper/xxxCG.tex} stops at the SoV pseudovacuum/covector construction and does
not yet state a TQ equation.

As recorded in \texttt{HANDOFF.md}, this notebook's cell outputs currently render as empty on disk
--- a wolfbook tool-capture glitch during the merge, confirmed unrelated to correctness by direct
kernel-state inspection (independently re-confirmed here via \texttt{wolfbook\_getKernelState}: the
transfer-matrix eigensystem, the coefficient extraction, and the Q-solve routines are all present as
live memoized kernel values, matching \texttt{Experiments/Baxter\_Solver\_L2.wb}'s numbers exactly).
Because of this, several worked examples below are sourced from the \textbf{predecessor}
\texttt{Experiments/Baxter\_Solver\_L2.wb} (which does display outputs) rather than from this
notebook's own currently-blank cell outputs; \texttt{HANDOFF.md} records that these were
spot-checked to match exactly during promotion. A manual ``Run All'' in the front end will
repopulate this notebook's own displayed outputs.

\paragraph{A note on notation.} The notebook itself uses literal Greek Unicode characters
($\lambda$, $\theta$, $\kappa$, $\chi$, $\tau$) as Wolfram symbol names -- e.g.\ the representation
labels are the Wolfram symbols whose names literally are $\lambda_1,\lambda_2$, the inhomogeneities
are $\theta_1,\theta_2$, the twist eigenvalues are $\kappa_1,\kappa_2$, the Newton-identity
coefficients of the twist are $\chi_1,\chi_2$, and the transfer-matrix eigenvalue function is named
$\tau$. Throughout this document these are rendered as their corresponding math-mode macros rather
than reproduced as literal glyphs; only Wolfram function calls with purely Latin/numeric arguments
(e.g.\ \texttt{TauEigensystem[1,1]}, \texttt{QSolve[2,1][2,0]}) are quoted verbatim in
\texttt{typewriter} font.

\section{The Gelfand--Tsetlin basis and single-site \texorpdfstring{$\mathfrak{gl}(2)$}{gl(2)} generators}
\label{sec:gt}

The single-site building blocks are ported verbatim from \texttt{Experiments/XXX\_su2.wb} and are
independent of the chain length $L$; every later construction in the notebook is built on top of
them. The Wolfram routine \texttt{ps[rep]} enumerates the Gelfand--Tsetlin (GT) patterns of a
rectangular $\mathfrak{gl}(N)$ irrep specified by a shape such as $\{\lambda,0\}$, memoized;
\texttt{dim[rep]} is simply the length of that list, i.e.\ the irrep dimension (memoized). For the
two-row $\mathfrak{gl}(2)$ shapes used throughout this notebook this gives the expected
$\dim V_\lambda = \lambda+1$: the notebook's own kernel state confirms $\dim\{1,0\}=2$,
$\dim\{2,0\}=3$, $\dim\{3,0\}=4$.

On top of this basis, \texttt{J0[k][pat]}, \texttt{Jp[k][pat]}, and \texttt{Jm[k][pat]} implement
the Cartan, raising, and lowering generator action on a single GT pattern \texttt{pat} at GT-row
level $k$, each returning a formal sum of basis vectors. The routine \texttt{JJ[k,l,rep]} assembles
these into the full matrix of the generator $E_{kl}$ (in the standard $\mathfrak{gl}(N)$ generator
convention) on \texttt{rep}'s GT basis: diagonal entries ($k=l$) and the entries adjacent to the
diagonal ($l=k\pm1$) are read directly from \texttt{Jp}/\texttt{Jm}, while all remaining entries are
built recursively from commutators of adjacent generators. There is no dedicated check cell for this
machinery in isolation; its correctness is implicit in every downstream check (most directly the
transfer-matrix commutativity check of \S\ref{sec:twosite} below, which would fail if the generator
matrices were wrong). As a representative evaluation, the notebook's kernel state gives
\begin{verbatim}
JJ[1, 1, {2, 0}]
(* {{2, 0, 0}, {0, 1, 0}, {0, 0, 0}} *)
\end{verbatim}
i.e.\ the diagonal weight (Cartan) values $2,1,0$ for the three $\lambda=2$ GT basis states. This
GT machinery underlies the phrase ``realized on the Gelfand--Tsetlin basis'' in
\texttt{paper/xxxCG.tex} \S2 (``Lax matrix and transfer matrix''), though the paper itself does not
spell out the GT formulas -- this is standard representation theory rather than a numbered paper
equation. All of the above is defined in Cell 3 of the notebook.

A thin wrapper, $E_{ij}\in\mathrm{End}(V_\lambda)$ (the Wolfram routine \texttt{Ee[$\lambda$][i,j]},
defined in Cell 4), simply calls \texttt{JJ[i,j,\{$\lambda$,0\}]} to give the spin-$\lambda/2$
representation matrix of the generator $E_{ij}$; this is exactly the object
$E_{ij}\in\mathrm{End}(V_\lambda)$ appearing in \texttt{paper/xxxCG.tex} \S2 immediately after the
Lax-matrix formula.

The single-site Lax matrix (Cell 5, Wolfram routine \texttt{L[$\lambda$][i,j][u]}) is then
\begin{equation}
  L^{[\lambda]}(u) \;=\; u\,\mathbbm{1} \;-\; h \sum_{i,j} E_{ji}\otimes e_{ij},
  \label{eq:Lax}
\end{equation}
matching \texttt{paper/xxxCG.tex} \S2; the $(i,j)$ auxiliary-space component implemented by the
routine is precisely the $e_{ij}$-coefficient of \eqref{eq:Lax},
\[
  L^{[\lambda]}_{ij}(u) = u\,\delta_{ij}\,\mathbbm{1} - h\,E_{ji}.
\]
Here $u$ is left unshifted at this stage; the inhomogeneity shift $u\mapsto u-\theta$ is applied by
the caller when a site is actually instantiated in the two-site chain (\S\ref{sec:twosite} below).

\section{The companion twist and the two-site transfer matrix}
\label{sec:twosite}

The companion-form twist matrix $G$ (Cell 6, with numeric eigenvalues $\kappa_1,\kappa_2$ fixed in
Cell 1) is
\begin{equation}
  G = \begin{pmatrix} \chi_1 & -\chi_2 \\ 1 & 0 \end{pmatrix}, \qquad
  \chi_1=\kappa_1+\kappa_2, \qquad \chi_2=\kappa_1\kappa_2,
  \label{eq:G}
\end{equation}
exactly as in \texttt{paper/xxxCG.tex} \S2.2 (``Companion twist and the fusion hierarchy''). In this
notebook $\kappa_2=1/\kappa_1$ is chosen so $\chi_2=1$. The construction is checked in
\texttt{Experiments/Baxter\_Solver\_L2.wb} (Cell 6), which verifies
$\{\operatorname{tr}G,\det G\}=\{\kappa_1+\kappa_2,\ \kappa_1\kappa_2\}$ after \texttt{Chop}; trace
and determinant are used rather than sorted eigenvalues because $\kappa_2=\overline{\kappa_1}$
exactly and numeric eigenvalue extraction introduces $\sim10^{-17}$ real-part noise that can flip a
naive sort's tie-break.

The two-site Lax operators (Cell 8, Wolfram routines \texttt{L1}, \texttt{L2}) act on
$V_{\lambda_1}\otimes V_{\lambda_2}$ by shifting each site's single-site Lax matrix by its own
inhomogeneity and tensoring with the identity on the other factor,
\[
  L_1(u) = L^{[\lambda_1]}(u-\theta_1)\otimes \mathbbm{1}_{\lambda_2}, \qquad
  L_2(u) = \mathbbm{1}_{\lambda_1}\otimes L^{[\lambda_2]}(u-\theta_2).
\]
Both $\lambda_1$ and $\lambda_2$ must be supplied to each routine even though only one representation
enters the shifted Lax matrix itself, because the size of the identity factor depends on the
\emph{other} site's representation. The two-site monodromy matrix (Cell 10, memoized routine
\texttt{T}) is the ordered product in auxiliary space,
\begin{equation}
  T(u) = L_1(u)\, L_2(u),
  \label{eq:monodromy}
\end{equation}
with both factors carrying the \emph{same} spectral parameter $u$ (each site's inhomogeneity shift
is already baked into $L_1$, $L_2$); this differs from the single-site quantum-determinant
construction, which shifts the second Lax evaluation by $h$ -- that shift is between two evaluations
of the \emph{same} site's Lax matrix and simply does not arise for two distinct physical sites, as
the notebook's own Cell 9 comment notes.

The fundamental, companion-twisted transfer matrix (Cell 11) is
\begin{equation}
  t_{1,1}(u) = \operatorname{tr}_a\!\bigl[G\,T(u)\bigr],
  \label{eq:t11}
\end{equation}
generalizing the single-site formula $t_{1,1}(u)=\operatorname{tr}_a[G\,L^{[\lambda]}(u-\theta)]$ of
\texttt{paper/xxxCG.tex} \S2 to the two-site monodromy \eqref{eq:monodromy}; as a matrix on
$V_{\lambda_1}\otimes V_{\lambda_2}$ it is a degree-2 polynomial in $u$. Commutativity of this
family at different spectral parameters -- the statement that its expansion coefficients are honest
commuting integrals of motion -- is checked in Cell 12 for $\lambda_1,\lambda_2\in\{1,\dots,4\}$,
designed to collapse $t_{1,1}(u)t_{1,1}(v)-t_{1,1}(v)t_{1,1}(u)$ to zero after simplification; the
check ran without error in the fresh-kernel top-to-bottom pass recorded in \texttt{HANDOFF.md} (its
own displayed output is currently blank, see the Overview).

A purely symbolic variant with a formally \emph{diagonal} twist (Cell 13, routine \texttt{t2diag})
replaces $G$ by two independent symbolic eigenvalues, written here as $\kappa(1),\kappa(2)$ to
emphasize that they are \textbf{distinct symbols} from the numeric plain twist eigenvalues
$\kappa_1,\kappa_2$ fixed in Cell 1 -- this variant is used only for a symbolic large-$u$
asymptotics cross-check and does not otherwise interact with the rest of the notebook. The
corresponding total (co-product) generator (Cell 14, routine \texttt{EE}),
\[
  E_{ij}^{\mathrm{tot}} = E_{ij}\otimes\mathbbm{1} + \mathbbm{1}\otimes E_{ij},
\]
satisfies the central-charge identity $E_{11}^{\mathrm{tot}}+E_{22}^{\mathrm{tot}}=(\lambda_1+\lambda_2)\,\mathbbm{1}$,
checked in Cell 15 over $\lambda_1,\lambda_2\in\{1,\dots,4\}$ (designed to collapse to zero; its
displayed output is currently blank for the same cosmetic reason). The predicted large-$u$
asymptotic expansion of the diagonal-twist transfer matrix up to and including the $u^1$ term (Cell
16, routine \texttt{tdiagasympt}) is
\[
  t^{\mathrm{diag}}_{1,1}(u) \;\sim\; \bigl(\kappa(1)+\kappa(2)\bigr)u^2\,\mathbbm{1}
  - u\bigl(\kappa(1)+\kappa(2)\bigr)(\theta_1+\theta_2)\,\mathbbm{1}
  - u\,h\bigl(\kappa(1)\,E_{11}^{\mathrm{tot}} + \kappa(2)\,E_{22}^{\mathrm{tot}}\bigr),
\]
and Cell 17 confirms that the $u^1$ coefficient of this expansion agrees with the direct
$t^{\mathrm{diag}}_{1,1}$ construction for $(\lambda_1,\lambda_2)=(1,1)$.

\section{Quantum determinant and the fusion hierarchy}
\label{sec:fusion}

The (untwisted) quantum determinant of the two-site monodromy matrix (Cell 18, routine
\texttt{qdetT}) is
\begin{equation}
  \mathrm{qdet}\,T(u) = T_{11}(u)\,T_{22}(u-h) - T_{21}(u)\,T_{12}(u-h),
  \label{eq:qdet}
\end{equation}
generalizing the single-site formula of \texttt{paper/xxxCG.tex} \S2.2,
$\mathrm{qdet}\,L^{[\lambda]}(u)=L^{[\lambda]}_{11}(u)L^{[\lambda]}_{22}(u-h)-L^{[\lambda]}_{21}(u)L^{[\lambda]}_{12}(u-h)$,
to the two-site monodromy. Cell 20 then builds the Hirota (T-system) fusion hierarchy extending
$t_{1,1}$ to general rectangular auxiliary representations $t_{a,s}(u)$, exactly as in the $L=1$
case: the boundary conditions $t_{0,s}\equiv t_{a,0}\equiv\mathbbm{1}$; the antisymmetric fusion
\begin{equation}
  t_{2,1}(u) = \chi_2\,\mathrm{qdet}\,T(u);
  \label{eq:t21}
\end{equation}
and the symmetric fusions via the Hirota recursion
\begin{equation}
  t_{1,s}(u) = t_{1,s-1}(u)\,t_{1,1}\bigl(u+(s-1)h\bigr) - t_{2,1}\bigl(u+(s-1)h\bigr)\,t_{1,s-2}(u),
  \qquad s>1.
  \label{eq:hirota}
\end{equation}
This is an exact generalization of \texttt{paper/xxxCG.tex} \S2.2, which states \eqref{eq:t21}--\eqref{eq:hirota}
for the single-site quantum determinant, to the two-site case.

\section{Separation of variables: the \texorpdfstring{$B$}{B}-operator and covector basis}
\label{sec:sov}

Two $\qth$-type normalization factors are defined in Cell 21: the single-site factor
$Q_{\theta,1}(u)=u-\theta_1$ and the two-site product $\qth(u)=(u-\theta_1)(u-\theta_2)$,
generalizing the single-site definition $\qth(u)=u-\theta$ of \texttt{paper/xxxCG.tex} \S2.3. Only
$\qth$ is actually used downstream (the $x_2$ construction below); $Q_{\theta,1}$ is defined but not
called anywhere else in this notebook, apparently kept for parity with the single-site file it was
ported from. Note also that $\qth$ is \textbf{redefined identically} in Section B of the notebook
(Cell 28, after its own \texttt{ClearAll}) -- harmless while the two definitions agree, but an edit
to one without the other would silently diverge.

The SoV $B$-operator for $L=2$ (Cell 23, routine \texttt{B2}) is simply the bare $(1,1)$ monodromy
element, with no extra twist or trace,
\[
  B_2(u) = T_{11}(u),
\]
matching the single-site convention $B^{[\lambda]}(u)=L^{[\lambda]}_{11}(u-\theta)$ of
\texttt{paper/xxxCG.tex} \S2.3; no additional shift is needed here since $\theta_1,\theta_2$ are
already built into $T$ via $L_1,L_2$.

The SoV pseudovacuum and covector basis (Cell 25) is built in two steps. The single-site
lowest-weight covector of $V_\lambda^*$'s GT basis (routine \texttt{xSingle}) tensors to give the
two-site pseudovacuum,
\[
  \langle x_2(0,0)| = \langle x^{[\lambda_1]}_0|\otimes\langle x^{[\lambda_2]}_0|,
\]
and the full covector basis is obtained (memoized routine \texttt{x2}) by acting with the fused
transfer matrices and normalizing,
\begin{equation}
  \langle x_2(s_1,s_2)| =
  \frac{\langle x_2(0,0)|\;t_{1,s_1}(\theta_1)\;t_{1,s_2}(\theta_2)}
       {\displaystyle\prod_{k=1}^{s_1}\qth\bigl(\theta_1+(k-1-\lambda_1)h\bigr)\;
        \prod_{k=1}^{s_2}\qth\bigl(\theta_2+(k-1-\lambda_2)h\bigr)}.
  \label{eq:x2}
\end{equation}
This generalizes the single-site formula of \texttt{paper/xxxCG.tex} \S2.3,
$\langle x_s^{[\lambda]}| = \langle x_0^{[\lambda]}|\,t_{1,s}(\theta)\big/\prod_{p=1}^{s}\qth(\theta-\lambda h+(p-1)h)$,
to $L=2$ via a tensor-product pseudovacuum with independent fusion actions and normalizations at
each site. Cell 26 verifies, for $s_1\in\{0,\dots,\lambda_1\}$, $s_2\in\{0,\dots,\lambda_2\}$,
$\lambda_1,\lambda_2\in\{1,2,3\}$, that \eqref{eq:x2} diagonalizes $B_2$ with the expected
eigenvalue,
\[
  \langle x_2(s_1,s_2)|\, B_2(u) = (u-\theta_1-hs_1)(u-\theta_2-hs_2)\,\langle x_2(s_1,s_2)|,
\]
the product of the two per-site eigenvalues $(u-\theta-sh)$ from \texttt{paper/xxxCG.tex}, reflecting
the tensor-product SoV structure; the check is designed to collapse to zero and ran without error in
the recorded full pass (its own displayed output is currently blank, cosmetic issue as above).

\section{The Baxter TQ equation (notebook Section B)}
\label{sec:baxter}

The routines in this final part of the notebook jointly solve the transfer-matrix eigenproblem
$t_{1,1}(u)\Psi=\tau(u)\Psi$ together with the Baxter TQ equation
\begin{equation}
  \kappa_1\,a(u)\,q(u+h) + \kappa_2\,\qth(u)\,q(u-h) = \tau(u)\,q(u),
  \label{eq:TQ}
\end{equation}
for the two-site chain, in two stages. First, \texttt{TauEigensystem} diagonalizes the
\emph{genuine} transfer matrix directly -- with no oversized Baxter-matrix proxy and no post-hoc
spurious-root filtering -- to obtain every $(M,n)$-indexed eigenvalue $\tau=\tau_0+J_n$. Then
\texttt{QSolve} solves for the monic degree-$M$ Q-function $q_n(u)$ as a linear nullspace problem at
the now-known $J_n$. This supersedes an earlier approach (still present only in
\texttt{Experiments/}) that diagonalized an oversized $(M+1)$-dimensional Baxter matrix directly and
could manufacture spurious extra roots. \texttt{paper/xxxCG.tex} does not yet state a Baxter/TQ
equation for either $L=1$ or $L=2$, so none of this section has a paper cross-reference; the full
derivation writeup and validation/regression suite live in
\texttt{Experiments/Baxter\_Solver\_L2.wb} rather than being duplicated into this Clean notebook.

\subsection*{Baxter-equation ingredients}

Cell 28 defines, locally so that Section B is self-contained: the Bethe/Baxter coefficient
\[
  a(u) = (u-\theta_1-h\lambda_1)(u-\theta_2-h\lambda_2);
\]
the $\qth$ factor $\qth(u)=(u-\theta_1)(u-\theta_2)$, an identical redefinition of the one from
\S\ref{sec:sov} (see the gotcha noted there); the $J$-independent part of the transfer-matrix
eigenvalue,
\[
  \tau_0(u) = (\kappa_1+\kappa_2)u^2 - (\theta_1+\theta_2)(\kappa_1+\kappa_2)u
    - u\,h\bigl(\kappa_1(\lambda_1+\lambda_2-M)+\kappa_2 M\bigr),
\]
which by construction has zero constant term, so that the full eigenvalue is $\tau=\tau_0+J$ with
$J$ the constant term; and a helper \texttt{qpoly[clist,uu]} that evaluates a polynomial given
ascending coefficients $\{c_0,\dots,c_M\}$ at a point, avoiding an indeterminate $0^0$. This is the
standard XXX Baxter-equation form, $a(u)=\prod_i(u-\theta_i-h\lambda_i)$,
$\qth(u)=\prod_i(u-\theta_i)$, generalized to two inhomogeneous sites with a companion twist of
eigenvalues $\kappa_1,\kappa_2$.

\subsection*{The Baxter matrix and transfer-matrix coefficients}

\texttt{BaxterMatrix} at fixed $(\lambda_1,\lambda_2,M)$ (Cell 29) builds the fixed linear operator, in the
monic-degree-$M$-polynomial coefficient basis, whose nullspace at a known $J$ gives the
Q-function's coefficients. For $k=0,\dots,M$ its $k$-th column is obtained by expanding
\begin{equation}
  O[u^k] = \kappa_1\,a(u)(u+h)^k + \kappa_2\,\qth(u)(u-h)^k - \tau_0(u)\,u^k
  \label{eq:baxtercol}
\end{equation}
as a polynomial in $u$; the routine hard-asserts (aborting on failure) that the degree-$(M{+}1)$ and
degree-$(M{+}2)$ coefficients of \eqref{eq:baxtercol} vanish identically for every column -- exactly
where a wrong $a$, $\qth$, $\tau_0$, or a wrong $\lambda_1+\lambda_2-M$ convention would first show
up -- and the surviving degree-$0,\dots,M$ coefficients become the columns of the resulting
$(M{+}1)\times(M{+}1)$ matrix. A smoke test in Cell 29,
\texttt{Dimensions /@ (BaxterMatrix[1,1,\#]\&/@\{0,1,2\})}, gives $\{\{1,1\},\{2,2\},\{3,3\}\}$,
confirmed identically in the predecessor \texttt{Experiments/Baxter\_Solver\_L2.wb} (Cells 4--5),
which additionally shows the cancelled high-degree coefficients numerically at $\sim10^{-29}$ (i.e.\
zero at the tracked 30-digit precision) for $M=0,1,2$ at $(\lambda_1,\lambda_2)=(1,1)$.

Separately, \texttt{t2Coeffs[$\lambda_1$,$\lambda_2$]} (Cell 30, memoized) extracts the coefficient
matrices $C_0,C_1,C_2$ of
\[
  t_{1,1}(u) = C_0 + C_1 u + C_2 u^2
\]
entrywise via \texttt{CoefficientList} in a dummy variable, reusable across every $(M,n)$ of that
representation. Its own check confirms $C_2=(\kappa_1+\kappa_2)\mathbbm{1}$ -- the transfer matrix's
leading-order, twist-independent asymptotics -- and is confirmed identically in
\texttt{Experiments/Baxter\_Solver\_L2.wb} (Cell 7), which returns
$\max|C_2-(\kappa_1+\kappa_2)\mathbbm{1}|\to 0\times10^{-29}$ for $(\lambda_1,\lambda_2)=(1,1)$.

The predicted dimension of the weight-$M$ subspace of $V_{\lambda_1}\otimes V_{\lambda_2}$ (Cell 31,
routine \texttt{dPred}) is the pure combinatorial count of $(m_1,m_2)$ with $m_1+m_2=M$ and
$0\le m_i\le\lambda_i$,
\[
  d_{\mathrm{pred}}(M,\lambda_1,\lambda_2) = \min(M,\lambda_1,\lambda_2,\lambda_1+\lambda_2-M) + 1,
\]
which is the exact multiplicity that \texttt{TauEigensystem} below must reproduce per $M$-sector.

\subsection*{Diagonalizing the genuine transfer matrix}

\texttt{TauEigensystem[$\lambda_1$,$\lambda_2$]} (Cell 31, memoized) diagonalizes the genuine
transfer matrix $t_{1,1}(u)$ directly. Since $C_0,C_1,C_2$ pairwise commute, it diagonalizes a
generic linear combination $C_0+rC_1$, trying $r\in\{7/3,\pi,e,11/5,13/4,\sqrt2\}$ in turn until the
combination's eigenvectors have full rank $d=(\lambda_1+1)(\lambda_2+1)$ with $d$ distinct
eigenvalues, guarding against an accidental degeneracy at a particular $r$. For each eigenvector $v$
it extracts
\[
  t_0 = \frac{v\cdot C_0 v}{v\cdot v}, \qquad t_1 = \frac{v\cdot C_1 v}{v\cdot v},
\]
and hard-aborts if either residual $\|C_0v-t_0v\|$ or $\|C_1v-t_1v\|$ exceeds $10^{-15}$; this
residual check \emph{is} the genuineness test, replacing an earlier post-hoc spurious-state filter.
The weight label $M$ is then read off algebraically from $t_1$ by inverting $\tau_0$'s $u^1$
coefficient,
\begin{equation}
  M = \frac{-(\theta_1+\theta_2)(\kappa_1+\kappa_2) - h\kappa_1(\lambda_1+\lambda_2) - t_1}
           {h(\kappa_2-\kappa_1)},
  \label{eq:Mextract}
\end{equation}
hard-asserted to lie within $10^{-10}$ of an integer, and $J=t_0$ directly (since $\tau_0$'s $u^0$
coefficient vanishes by construction). States within each $M$-sector are sorted ascending by
$(\operatorname{Re}J,\operatorname{Im}J)$ to assign the index $n$, and each sector's count is
hard-asserted to equal $d_{\mathrm{pred}}(M,\lambda_1,\lambda_2)$ exactly. The result is an
association keyed by $(M,n)$ giving the pair $(J_n,\Psi_{M,n})$, with $\Psi_{M,n}$ normalized so its
last component is $1$.

These in-function hard asserts (rank/degeneracy resolution, residual below $10^{-15}$, integer $M$
within $10^{-10}$, exact per-sector counts) constitute the check: a genuine failure aborts rather
than silently returning wrong data. Independently, \texttt{Experiments/Baxter\_Solver\_L2.wb}
(Cells 10--11) re-checks $t_{1,1}(u)\Psi-\tau(u)\Psi$ from scratch (not via the $C_0,C_1,C_2$
extraction used internally) at five sample points $u\in\{0,1,2,-1,5\}$ for
$(\lambda_1,\lambda_2)=(1,1)$ and $(2,1)$, with maximum residual $\sim10^{-26}$--$10^{-27}$ over all
genuine states and genuine-state counts per $M$ matching $d_{\mathrm{pred}}$ exactly in both cases.
This notebook's own live kernel state already carries \texttt{TauEigensystem[1,1]} and
\texttt{TauEigensystem[2,1]} with $J$-values numerically identical to
\texttt{Baxter\_Solver\_L2.wb}'s, e.g.
\begin{verbatim}
KeySort[TauEigensystem[1, 1]][[All, "J"]] // N
(* <|{0,0} -> 0.5664055180884988+1.3769645669711723 I,
     {1,0} -> -0.778123966451158+0. I,
     {1,1} -> 1.1900552523337025+0. I,
     {2,0} -> 0.5664055180884988-1.3769645669711723 I|> *)
\end{verbatim}

Bounds-checked, on-demand accessors over the shared \texttt{TauEigensystem} cache (Cell 32) then
give the eigenvector, $\texttt{Psi}(M,n)=\Psi_{M,n}$, and the full eigenvalue polynomial,
$\tau_{M,n}(u)=\tau_0(u)+J_n$, respectively. Both require $M$ an integer in $0,\dots,\lambda_1+\lambda_2$
and $n$ an integer in $0,\dots,M$, else a bounds message and \texttt{\$Failed}; if $n$ is in this
gross range but exceeds the actual weight multiplicity $d_{\mathrm{pred}}(M,\lambda_1,\lambda_2)-1$,
a spurious-state message fires and a \texttt{Missing} value is returned -- the only remaining
``spurious'' outcome, since by construction every state \texttt{TauEigensystem} does produce is
genuine. A deliberate smoke test in Cell 32 calls \texttt{Psi[2,1][2,2]} (which passes the gross
bound $n\le M=2$ but has $d_{\mathrm{pred}}(2,2,1)=2$, so only $n=0,1$ actually exist), and the
spurious-state message and missing result are confirmed there and cross-checked identically in
\texttt{Experiments/Baxter\_Solver\_L2.wb} (Cell 9).

\subsection*{Solving for the Q-function}

\texttt{QSolve[$\lambda_1$,$\lambda_2$][M,n]} (Cell 33, memoized) solves \eqref{eq:TQ} as a
\emph{linear} problem rather than an eigenvalue problem, since $J_n$ (hence $\tau$) is already known
from \texttt{TauEigensystem}: it finds the monic degree-$M$ polynomial $q_n$ satisfying
\[
  \bigl(\mathrm{BaxterMatrix}(\lambda_1,\lambda_2,M) - J_n\mathbbm{1}\bigr)\, q_n = 0,
\]
a matrix that is singular by construction. The routine's own comment records an important precision
gotcha: the would-be-zero entries of this matrix come out as \emph{precision-tracked} zeros (exact
value $0$, but carrying a finite Accuracy/Precision tag, e.g.\ printed as $0\times10^{-26}$), and a
bare nullspace computation treats such an entry as unknown rather than confirmed-zero -- this
specifically caused a failure in the $M=0$ sector (a $1\times1$ matrix whose single entry is such a
precision-tracked zero) under either the default tolerance or an explicit tolerance of $10^{-10}$
alone. The fix used inside \texttt{QSolve} is to \texttt{Chop} the matrix at a threshold of
$10^{-12}$ -- comfortably above the $\sim10^{-26}$--$10^{-29}$ noise floor but far below any genuine
$O(1)$ entry -- \emph{before} computing the nullspace, converting those entries to exact zeros; an
explicit tolerance of $10^{-10}$ is kept in addition for headroom against genuinely small but
nonzero near-null directions that chopping alone would not catch in general. This is confirmed as an
explicit regression guard in \texttt{Experiments/Baxter\_Solver\_L2.wb} (Cell 13): for
$(\lambda_1,\lambda_2,M,n)=(2,1,0,0)$, a bare nullspace computation and one with an explicit
tolerance of $10^{-10}$ alone both return an empty nullspace (dimension $0$), while chopping together
with the tolerance (as used inside \texttt{QSolve}) correctly returns dimension $1$. As a
representative, regression-critical example, reproduced identically in this notebook's own kernel
state,
\begin{verbatim}
QSolve[2, 1][2, 0]
(* {1.03464899078135229594035121+0.381869568380946283899575941 I,
    -2.37022278949001677248172619-0.20048970475590097847093779 I,
    1.00000000000000000000000000+0.\times10^-27 I} *)
\end{verbatim}
gives the ascending coefficients of the monic degree-2 Q-polynomial for that sector.

Finally, \texttt{Q1[$\lambda_1$,$\lambda_2$][M,n][u]} (Cell 34) assembles the actual Baxter
Q-function value from the same cache,
\[
  Q_{M,n}(u) = \kappa_1^{u/h}\, q_n(u),
\]
mutually consistent with $\tau_{M,n}$ by construction since both trace back to the same
\texttt{TauEigensystem} entry. The full end-to-end TQ-equation residual,
\[
  a(u)\,Q(u+h) - \tau(u)\,Q(u) + \qth(u)\,Q(u-h) = 0,
\]
assembled entirely from \texttt{Q1} and $\tau$, is checked in
\texttt{Experiments/Baxter\_Solver\_L2.wb} (Cells 15--16) at sample points avoiding the roots of $a$
and $\qth$, for every $(M,n)$ in both $(\lambda_1,\lambda_2)=(1,1)$ (maximum residual
$\sim10^{-21}$--$10^{-26}$) and $(2,1)$ (maximum residual $\sim10^{-20}$--$10^{-26}$, with the
historically-critical case $(M,n)=(2,0)$ individually flagged and confirmed). Per \texttt{HANDOFF.md},
this notebook's own \texttt{Q1}/\texttt{QSolve} outputs were spot-checked against
\texttt{Baxter\_Solver\_L2.wb}'s and match exactly. As a representative example, reproduced
identically in this notebook's kernel state,
\begin{verbatim}
Q1[1, 1][1, 0][2]
(* -0.709231010958016926537225142+1.09596916251277206706752728 I *)
\end{verbatim}
together with the corresponding transfer-matrix eigenvalue $\tau_{1,0}^{[1,1]}(2)\approx-0.02291660900$
to ten digits (obtained from $\tau_0(u)+J_{1,0}$ evaluated numerically), consistent with
\eqref{eq:TQ} to the tracked precision.

\section*{Known gotchas / non-obvious conventions}

\begin{itemize}
\item \textbf{Cell-output display glitch (cosmetic only).} This notebook's cell outputs currently
  render empty on disk (confirmed via direct inspection of every cell). This is a documented,
  understood tool-capture artifact from the merge, not a correctness problem -- direct kernel-state
  inspection shows all expected memoized values present and numerically matching the predecessor
  notebooks. A manual ``Run All'' will repopulate the displayed outputs.
\item \textbf{Symbolic $\kappa(1)$/$\kappa(2)$ vs.\ the numeric plain variables $\kappa_1$/$\kappa_2$.}
  The diagonal-twist transfer matrix and its asymptotic-expansion cross-check (\S\ref{sec:twosite})
  use a formal, purely symbolic indexed twist $\kappa(i)$, unrelated to the numeric plain twist
  eigenvalues $\kappa_1,\kappa_2$ set once at the top of the notebook and used everywhere else
  (including all of Section B). They only appear together, symbolically, in the one asymptotics
  cross-check cell.
\item \textbf{$\qth$ is defined twice with the same formula.} Both the SoV section
  (\S\ref{sec:sov}) and the Baxter-solver section (\S\ref{sec:baxter}) define
  $\qth(u)=(u-\theta_1)(u-\theta_2)$ independently (the latter after its own \texttt{ClearAll}).
  Currently harmless since the two agree, but a future edit to one without the other would silently
  diverge. The single-site factor $Q_{\theta,1}(u)=u-\theta_1$ from \S\ref{sec:sov} is defined but
  never used downstream in this notebook.
\item \textbf{Memoization caches stale results.} Several routines throughout the notebook
  (\texttt{ps}, \texttt{dim}, the GT generator routines, \texttt{T}, \texttt{x2}, \texttt{t2Coeffs},
  \texttt{TauEigensystem}, \texttt{QSolve}) are self-memoizing. Per the project-wide convention,
  after fixing any upstream definition, re-run the defining cell and everything downstream that used
  the stale cache -- a partial re-run will not clear old memoized values.
\item \textbf{Precision-tracked exact zeros before a nullspace computation.} See the \texttt{QSolve}
  discussion in \S\ref{sec:baxter}: always chop a numeric matrix before computing its nullspace when
  the singularity is expected by construction, rather than relying on an explicit tolerance option
  alone.
\item \textbf{High-precision twist, $h=1$.} The twist eigenvalue $\kappa_1$ is set to a 30-digit
  numeric value on the unit circle, $\kappa_1=e^{i\zeta(3)}$ (with $\kappa_2=1/\kappa_1=\overline{\kappa_1}$),
  and $h=1$. Numeric checks throughout rely on this tracked precision -- residuals reported at
  $10^{-20}$ to $10^{-29}$ are relative to this roughly 30-digit budget, not to machine precision --
  so lower-precision numeric twist values should not be substituted without re-checking the hard
  tolerance thresholds ($10^{-15}$, $10^{-10}$, $10^{-18}$, etc.) scattered through Section B.
\end{itemize}

\end{document}
